\documentclass{article}
\usepackage[T1]{fontenc}
\usepackage{lmodern}
\usepackage{graphicx}
\usepackage{amsmath,amssymb}
\usepackage[a4paper,margin=2cm]{geometry}
\usepackage[absolute,overlay]{textpos}
\setlength{\TPHorizModule}{1mm}
\setlength{\TPVertModule}{1mm}
\usepackage[indonesian]{babel}
\usepackage{hyperref}
\setlength{\emergencystretch}{2em}
% unit: Halaman 27

\begin{document}

% === Halaman 27 (llm) ===

\section*{Keystream Generator}

\begin{itemize}
    \item \textit{Keystream generator} diimplementasikan sebagai prosedur yang sama baik di sisi pengirim maupun di sisi penerima pesan.
    \item \textit{Keystream generator} dapat membangkitkan \textit{keystream} dalam satuan bit, atau \textit{byte}, atau blok bit.
    \item Bit-bit \textit{keystream} adalah bit-bit acak, namun bukan \textit{true random}, tetapi \textit{pseudo-random}, karena bit-bit \textit{keystream} dapat dibangkitkan ulang baik di sisi pengirim maupun di sisi penerima pesan asalkan umpan (\textit{seed}) U yang digunakan sama.
    \item Bit-bit \textit{pseudo-random} memiliki periode (dapat terulang kembali setelah beberapa kali iterasi)
\end{itemize}

\end{document}
